\begin{theorem}
\label{theorem-universally-exact-criteria}
Let
$$
0 \to M_1 \xrightarrow{f_1} M_2 \xrightarrow{f_2} M_3 \to 0
$$
be an exact sequence of $R$-modules.  The following are equivalent:
\begin{enumerate}
\item The sequence $0 \to M_1 \to M_2 \to M_3
\to 0$ is universally exact.
\item For every finitely presented $R$-module $Q$, the sequence
$$
0 \to M_1 \otimes_R Q \to M_2 \otimes_R Q \to
M_3 \otimes_R Q \to 0
$$
is exact.
\item Given elements $x_i \in M_1$ $(i = 1, \ldots, n)$, $y_j \in M_2$ $(j = 1,
\ldots, m)$, and $a_{ij} \in R$ $(i = 1, \ldots, n, j = 1, \ldots, m)$ such that
for all $i$
$$
f_1(x_i) = \sum\nolimits_j a_{ij} y_j,
$$
there exists $z_j \in M_1$ $(j =1, \ldots, m)$ such that for all $i$,
$$
x_i = \sum\nolimits_j a_{ij} z_j .
$$
\item Given a commutative diagram of $R$-module maps
$$
\xymatrix{
R^n \ar[r] \ar[d] &  R^m \ar[d] \\
M_1 \ar[r]^{f_1}        &  M_2
}
$$
where $m$ and $n$ are integers, there exists a map $R^m \to M_1$ making
the top triangle commute.
\item For every finitely presented $R$-module $P$, the $R$-module
map $\Hom_R(P, M_2) \to \Hom_R(P, M_3)$ is surjective.
\item The sequence $0 \to M_1 \to M_2 \to M_3
\to 0$ is the colimit of a directed system of split exact sequences of
the form
$$
0 \to M_{1} \to M_{2, i} \to M_{3, i} \to 0
$$
where the $M_{3, i}$ are finitely presented.
\end{enumerate}
\end{theorem}

\begin{proof}
Obviously (1) implies (2).

\medskip\noindent
Next we show (2) implies (3).  Let $f_1(x_i) = \sum_j a_{ij} y_j$ be relations
as in (3).  Let $(d_j)$ be a basis for $R^m$, $(e_i)$ a basis for $R^n$, and
$R^m \to R^n$ the map given by $d_j \mapsto \sum_i a_{ij} e_i$.
Let $Q$ be the cokernel of $R^m \to R^n$.  Then tensoring
$R^m \to R^n \to Q \to 0$ by the map $f_1: M_1 \to M_2$, we get a
commutative diagram
$$
\xymatrix{
M_1^{\oplus m} \ar[r] \ar[d] & M_1^{\oplus n} \ar[r] \ar[d] & M_1 \otimes_R Q
\ar[r] \ar[d] & 0 \\
M_2^{\oplus m} \ar[r] & M_2^{\oplus n} \ar[r] & M_2 \otimes_R Q \ar[r] & 0
}
$$
where $M_1^{\oplus m} \to M_1^{\oplus n}$ is given by
$$
(z_1, \ldots, z_m) \mapsto
(\sum\nolimits_j a_{1j} z_j, \ldots, \sum\nolimits_j a_{nj} z_j),
$$
and $M_2^{\oplus m} \to M_2^{\oplus n}$ is given similarly.  We want to
show $x = (x_1, \ldots, x_n) \in M_1^{\oplus n}$ is in the image of $M_1^{\oplus
m} \to M_1^{\oplus n}$.  By (2) the map $M_1 \otimes Q \to M_2
\otimes Q$ is injective, hence by exactness of the top row it is enough to show
$x$ maps to $0$ in $M_2 \otimes Q$, and so by exactness of the bottom row it is
enough to show the image of $x$ in $M_2^{\oplus n}$ is in the image of
$M_2^{\oplus m} \to M_2^{\oplus n}$.  This is true by assumption.

\medskip\noindent
Condition (4) is just a translation of (3) into diagram form.

\medskip\noindent
Next we show (4) implies (5). Let $\varphi : P \to M_3$ be a map from a
finitely presented $R$-module $P$.  We must show that $\varphi$ lifts to a map
$P \to M_2$.  Choose a presentation of $P$,
$$
R^n \xrightarrow{g_1} R^m \xrightarrow{g_2} P \to 0.
$$
Using freeness of $R^n$ and $R^m$, we can construct $h_2: R^m \to M_2$
and then $h_1: R^n \to M_1$ such that the following diagram commutes
$$
\xymatrix{
         & R^n \ar[r]^{g_1} \ar[d]^{h_1} & R^m \ar[r]^{g_2} \ar[d]^{h_2} & P
\ar[r] \ar[d]^{\varphi} & 0 \\
0 \ar[r] & M_1 \ar[r]^{f_1} & M_2 \ar[r]^{f_2} & M_3 \ar[r] & 0 .
}
$$
By (4) there is a map $k_1: R^m \to M_1$ such that $k_1 \circ g_1 =
h_1$.  Now define $h'_2: R^m \to M_2$ by $h_2' = h_2 - f_1 \circ k_1$.
Then
$$
h'_2 \circ g_1 = h_2 \circ g_1 - f_1 \circ k_1 \circ g_1 =
h_2 \circ g_1 - f_1 \circ h_1 = 0 .
$$
Hence by passing to the quotient $h'_2$ defines a map $\varphi': P \to
M_2$ such that $\varphi' \circ g_2 = h_2'$.  In a diagram, we have
$$
\xymatrix{
R^m \ar[r]^{g_2} \ar[d]_{h'_2} & P \ar[d]^{\varphi} \ar[dl]_{\varphi'} \\
M_2 \ar[r]^{f_2} & M_3.
}
$$
where the top triangle commutes.  We claim that $\varphi'$ is the desired lift,
i.e.\ that $f_2 \circ \varphi' = \varphi$.  From the definitions we have
$$
f_2 \circ \varphi' \circ g_2 = f_2 \circ h'_2 =
f_2 \circ h_2 - f_2 \circ f_1 \circ k_1 = f_2 \circ h_2 =
\varphi \circ g_2.
$$
Since $g_2$ is surjective, this finishes the proof.

\medskip\noindent
Now we show (5) implies (6). Write $M_{3}$ as the colimit of a directed system
of finitely presented modules $M_{3, i}$, see
Lemma \ref{lemma-module-colimit-fp}. Let $M_{2, i}$ be the fiber product
of $M_{3, i}$ and $M_{2}$ over $M_{3}$---by definition this is the submodule of
$M_2 \times M_{3, i}$ consisting of elements whose two projections onto $M_3$
are equal. Let $M_{1, i}$ be the kernel of the projection
$M_{2, i} \to M_{3, i}$. Then we have a directed system of exact sequences
$$
0 \to M_{1, i} \to M_{2, i} \to M_{3, i} \to 0,
$$
and for each $i$ a map of exact sequences
$$
\xymatrix{
0  \ar[r] & M_{1, i} \ar[d] \ar[r]  &  M_{2, i}  \ar[r] \ar[d] & M_{3, i} \ar[d]
\ar[r] & 0 \\
0  \ar[r] & M_{1} \ar[r]  &  M_{2}  \ar[r] & M_{3} \ar[r] & 0
}
$$
compatible with the directed system.  From the definition of the fiber product
$M_{2, i}$, it follows that the map $M_{1, i} \to M_1$ is an isomorphism.
 By (5) there is a map $M_{3, i} \to M_{2}$ lifting $M_{3, i} \to
M_3$, and by the universal property of the fiber product this gives rise to a
section of $M_{2, i} \to M_{3, i}$.  Hence the sequences
$$
0 \to M_{1, i} \to M_{2, i} \to M_{3, i} \to 0
$$
split.  Passing to the colimit, we have a commutative diagram
$$
\xymatrix{
0  \ar[r] & \colim M_{1, i} \ar[d]^{\cong} \ar[r]  &  \colim M_{2, i}
 \ar[r]
\ar[d] & \colim M_{3, i} \ar[d]^{\cong} \ar[r] & 0 \\
0  \ar[r] & M_{1} \ar[r]  &  M_{2}  \ar[r] & M_{3} \ar[r] & 0
}
$$
with exact rows and outer vertical maps isomorphisms.  Hence $\colim
M_{2, i}
\to M_2$ is also an isomorphism and (6) holds.

\medskip\noindent
Condition (6) implies (1) by
Example \ref{example-universally-exact} (2).
\end{proof}